\documentclass[11pt]{article}
\usepackage[margin=1in]{geometry}
\usepackage{amsmath,amssymb,amsthm,mathtools,booktabs,microtype}
\usepackage[T1]{fontenc}
\usepackage{lmodern}
\usepackage{xurl}
\usepackage{hyperref}
\hypersetup{hidelinks,pdftitle={A counterexample to Kourovka Notebook Problem 16.45},pdfauthor={Alec Kriebel}}
\newtheorem{theorem}{Theorem}
\newtheorem{proposition}[theorem]{Proposition}
\newtheorem{lemma}[theorem]{Lemma}
\newtheorem{corollary}[theorem]{Corollary}
\theoremstyle{remark}\newtheorem{remark}[theorem]{Remark}
\newcommand{\F}{\mathbf F}
\newcommand{\mup}{\mu'}
\newcommand{\bfth}{b_{\mathrm f}}
\newcommand{\GL}{\operatorname{GL}}
\newcommand{\SL}{\operatorname{SL}}
\newcommand{\Sym}{\operatorname{Sym}}
\newcommand{\Core}{\operatorname{Core}}
\newcommand{\rankcap}{c_{\cap}}
\title{A counterexample to Kourovka Notebook Problem 16.45}
\author{Alec Kriebel\\\small Independent researcher\\\small\href{https://orcid.org/0009-0001-9320-500X}{ORCID: 0009-0001-9320-500X}}
\date{23 September 2026\\\small Version 1.0.1; unrefereed preprint}
\begin{document}
\maketitle
\begin{abstract}
Let $b(G)$ be the largest size of an inclusion-minimal base, over all
permutation representations of a finite group $G$, with bases taken for the
permutation image. Let $\mup(G)$ be the largest size of an independent subset
of $G$, not necessarily generating $G$. We construct an explicitly specified
affine group of order $100920$ for which
\[
                 b(G)=3<4=\mup(G).
\]
The upper bound covers faithful and nonfaithful actions and does not depend on
enumerating the subgroups of the affine group. Exact, independently implemented
Python and C++ verifiers certify the small matrix-group inputs and the explicit
witnesses. The obstruction is a subgroup of order two that is central in the complement
but not normal in the ambient affine group.
\end{abstract}

\section{Introduction and conventions}
Cameron's Problem 16.45 in the \emph{Kourovka Notebook} asks whether two
invariants of every finite group coincide: the largest size of an
inclusion-minimal permutation base, and the largest size of an independent
subset. The inequality $b(G)\le\mup(G)$ is known. We give a negative answer
by constructing an affine group $G=\F_{29}^2\rtimes H$, where
$H\cong\SL_2(5)$ acts through explicitly given matrices, and proving
\[
                 \bfth(G)=b(G)=3<4=\mup(G).
\]
Here $\bfth$ restricts the base invariant to faithful actions. No
minimum-order assertion is made.

The distinction between faithful and arbitrary actions is essential to the
proof: for the complement, $\bfth(H)=2$ whereas $b(H)=3$. Every subgroup of
$G$ not containing its translation subgroup has faithful base invariant at
most two. Restricting an irredundant intersection family to such a subgroup
bounds faithful bases in $G$; the normal-subgroup structure then covers all
nonfaithful actions. The upper bound is structural and requires no enumeration
of the subgroup lattice of the order-$100920$ affine group.

\subsection{Definitions and relation to earlier work}
A subset $S$ of a finite group $X$ is \emph{independent} if
$x\notin\langle S\setminus\{x\}\rangle$ for each $x\in S$; the empty subset is
allowed. Write $\mup(X)$ for its maximum possible size. For an action
$\rho:X\to\Sym(\Omega)$, a base has pointwise stabilizer $\ker\rho$ in $X$.
Minimal means inclusion-minimal. Write $b(X)$ for the maximum size of a minimal
base over all actions, and $\bfth(X)$ for the corresponding maximum restricted
to faithful actions. Even if $\Omega$ is infinite, a base contains a finite
subbase: choose a point moved by each nonidentity element of the finite
permutation image. Consequently every inclusion-minimal base is finite.

These conventions give $b(1)=\bfth(1)=\mup(1)=0$. Cameron's 2010 exposition
uses $b_2$ for our $b$ and $\mu^*$ for our $\mup$; his 2014 preprint uses $b_2$ and $\mu'$, and his
2024 published account uses $b_2$ and $\mu^\uparrow$, respectively
\cite{blog,paper,cameron2024}. These sources allow nonfaithful representations.
The foundational subgroup and Boolean meet-semilattice formulations below
are due to Cameron \cite{paper,cameron2024}; we include proofs to fix the
normality and action-kernel conventions used in our construction.

The September 2026 notebook lists Problem 16.45 without an attached solution
\cite{notebook}. A focused search of the notebook, its update, and the cited
primary literature found no earlier resolution. This reports the scope of
the search, not a guarantee of bibliographic priority.

\subsection{Subgroup formulations}
\begin{proposition}[Subgroup formulation]
For a finite group $X$, $b(X)$ is the largest $t$ for which subgroups
$L_1,\ldots,L_t$ satisfy
\[
 N=\bigcap_{i=1}^t L_i\lhd X,
 \qquad
 \bigcap_{j\ne i}L_j>N\quad(1\le i\le t).
\]
An empty intersection is $X$. The same description with $N=1$ gives
$\bfth(X)$, and
\begin{equation}\label{eq:quotients}
 b(X)=\max_{N\lhd X}\bfth(X/N).
\end{equation}
\end{proposition}
\begin{proof}
The point stabilizers of a minimal base have precisely these properties, with
$N$ the action kernel. Conversely, let such subgroups be given and let $X$ act
on the disjoint union of the left coset spaces $X/L_i$. Its kernel is
\[
 \bigcap_i\Core_X(L_i)
 =\Core_X\!\left(\bigcap_iL_i\right)=\Core_X(N)=N.
\]
Choose the identity coset in each component. Their common stabilizer is $N$,
and deleting any one enlarges it, so these points are a minimal base. This
argument uses cores only to compute the kernel; it does not replace the
irredundant family by a family of cores. Formula \eqref{eq:quotients} follows
by factoring an action through its kernel and, conversely, pulling back a
faithful quotient action. The empty family handles the trivial group and
trivial actions consistently.
\end{proof}

Call a subgroup family \emph{meet-irredundant} if every member is essential to
its total intersection, without any normality requirement. Let $\rankcap(X)$
be its largest size.

\begin{lemma}\label{lem:breadth}
For every finite group $X$, $\rankcap(X)=\mup(X)$. In particular,
$b(X)\le\mup(X)$.
\end{lemma}
\begin{proof}
For a meet-irredundant family $L_1,\ldots,L_t$, choose
\[
 x_i\in\left(\bigcap_{j\ne i}L_j\right)\setminus L_i.
\]
Every $x_j$ with $j\ne i$ lies in $L_i$, whereas $x_i$ does not. Therefore the
$x_i$ are distinct and independent. Conversely, from an independent set
$\{x_1,\ldots,x_t\}$ take $L_i=\langle x_j:j\ne i\rangle$. The element $x_i$
witnesses the essentiality of $L_i$. Neither direction asserts that the total
intersection is normal.
\end{proof}

\begin{lemma}[Restriction]\label{lem:restriction}
If $L_1,\ldots,L_t$ is meet-irredundant in $X$, then, for each $i$,
\[
 t-1\le\mup(L_i).
\]
If its total intersection is $1$, the stronger bound
$t-1\le\bfth(L_i)$ holds.
\end{lemma}
\begin{proof}
Inside $L_i$, the family $L_i\cap L_j$ for $j\ne i$ remains meet-irredundant:
the original witness for $L_j$ lies in $L_i$ and witnesses this assertion.
Its total intersection equals the original total intersection. Apply
Lemma~\ref{lem:breadth}, or the faithful subgroup formulation when that
intersection is $1$.
\end{proof}

For completeness, an irredundant family gives a top-preserving Boolean
meet-semilattice embedding
\[
 A\longmapsto\bigcap_{i\notin A}L_i\qquad(A\subseteq\{1,\ldots,t\}).
\]
The witness $x_i$ distinguishes images of sets differing at $i$; intersection
corresponds to intersection of index sets. The bottom is $\bigcap_iL_i$.
There is no assertion that joins are preserved, or that this bottom is the
identity. Thus the proof retains the precise semilattice convention in the
notebook.

\section{The complement and its two different base invariants}
All matrices in the construction below are over $\F_{29}$. Put
\[
 A=\begin{pmatrix}0&28\\1&0\end{pmatrix},\qquad
 B=\begin{pmatrix}2&7\\12&28\end{pmatrix},\qquad
 H=\langle A,B\rangle\le\SL_2(29),\qquad Z=\{I,-I\}.
\]
Direct multiplication gives
\begin{equation}\label{eq:triangle}
 A^2=B^3=(AB)^5=-I.
\end{equation}
The standard spherical triangle presentation
\[
 A_5=\langle a,b\mid a^2=b^3=(ab)^5=1\rangle
\]
and simplicity of $A_5$ show that $H/Z\cong A_5$: the projective image is a
nontrivial quotient of this presentation, since $A\notin Z$. Consequently
$|H|=120$. The explicit finite isomorphism described next also identifies $H$ with
$\SL_2(5)$; the structural proof below uses only the quotient $H/Z\cong A_5$.

There is also a wholly finite certification of these facts, independent of
the triangle presentation. The supplied verifiers generate exactly 120
matrices and construct an explicit isomorphism from $\SL_2(5)$ by mapping
\[
 \begin{pmatrix}0&4\\1&0\end{pmatrix}\longmapsto A,
 \qquad
 \begin{pmatrix}0&4\\1&1\end{pmatrix}\longmapsto B.
\]
The Python verifier additionally constructs the quotient map onto $A_5$ by
mapping $A$ to $(1\,2)(3\,4)$ and $B$ to $(1\,3\,5)$, and checks that its
kernel is exactly $Z$. Thus the small-group identifications can alternatively
be read as explicit finite certificates, not as catalogue lookups.

We record the elementary group facts needed below. Every proper subgroup of
$A_5$ lies in a subgroup isomorphic to $A_4$, $D_{10}$, or $S_3$. Here $D_{10}$
denotes a group of order 10. One way to see this is to use the natural action
on five letters. An intransitive subgroup either fixes a letter or has orbit
sizes $2+3$. A proper transitive subgroup has order $5,10,15,20$, or $30$.
Orders 15 and 20 would give a normal Sylow 5-subgroup, whose normalizer in
$A_5$ has order 10. Order 30 is impossible: an index-two quotient is trivial
on every 3-cycle, and 3-cycles generate $A_5$. The remaining cases lie in a
Sylow-5 normalizer. This normalizer has order 10 because the order-20
normalizer in $S_5$ contains odd permutations.

Each of $A_4,D_{10},S_3$ has independence number 2. For $A_4$, a generating
set contains a 3-cycle and some element outside its cyclic subgroup; these
two generate $A_4$, since there is no subgroup of order 6. Its proper
subgroups have independence number at most 2. For a dihedral group of order
$2q$, $q$ prime, a generating set contains either a nontrivial rotation and
a reflection, or two distinct reflections, and two such elements suffice.
It follows, by deleting one element from any independent set in $A_5$, that
\begin{equation}\label{eq:A5rank}
                        \mup(A_5)\le3.
\end{equation}
The natural degree-five action has a minimal base of size 3, so equality holds.
Simplicity, used above and below, also has an elementary check: the nonidentity
conjugacy-class sizes are $15,20,12,12$, and no proper nontrivial union of
classes together with the identity has order dividing 60.

\begin{lemma}\label{lem:H}
The complement satisfies
\[
        \mup(H)=b(H)=3,\qquad \bfth(H)=2.
\]
\end{lemma}
\begin{proof}
The only involution in $\SL_2(29)$ is $-I$: a matrix squaring to $I$ is
diagonalizable with eigenvalues $\pm1$, and determinant one forces both
signs to agree. Therefore every even-order subgroup of $H$ contains $Z$.
Every odd-order subgroup injects into $A_5$ and, by the preceding subgroup
description, is trivial or cyclic of order 3 or 5.

Let $S$ be independent in $H$. If its images in $H/Z$, as an indexed family,
are independent, then $|S|\le3$. Otherwise some $s\in S$ lies in $KZ$, where
$K=\langle S\setminus\{s\}\rangle$, but does not lie in $K$. Thus $Z$ is not
contained in $K$, so $K$ has odd order and is trivial or cyclic of prime
order. Since $S\setminus\{s\}$ is independent, $|S|\le2$. This proves
$\mup(H)\le3$.

For the faithful invariant, a subgroup family with intersection $1$ must
have an odd-order member, since all even-order members contain $Z$. If this
member is trivial, a minimal such family has size 1. If it has prime order,
some other member intersects it trivially, so a minimal family has size at
most 2. Nontrivial subgroups of orders 3 and 5 give a faithful family of size
2, proving $\bfth(H)=2$.

For explicit witnesses to the remaining lower bounds, let
\begin{equation}\label{eq:R}
 R_1=\begin{pmatrix}0&28\\1&0\end{pmatrix},\qquad
 R_2=\begin{pmatrix}12&24\\0&17\end{pmatrix},\qquad
 R_3=\begin{pmatrix}25&3\\4&4\end{pmatrix}.
\end{equation}
They belong to $H$, since $R_1=A$, $R_2=BAB^2$ and
$R_3=AR_2(AB)^2$. Direct multiplication gives $R_i^2=-I$ and
\[
 |\langle R_2,R_3\rangle|=8,\qquad
 |\langle R_1,R_3\rangle|=12,\qquad
 |\langle R_1,R_2\rangle|=20.
\]
In the same pair order, the products $R_2R_3$, $R_1R_3$, and $R_1R_2$
have orders $4,3,10$. In $H/Z$, each pair consists of two involutions whose
product has order $2,3,5$, respectively. Thus the pairs generate dihedral
groups of orders $4,6,10$ in $H/Z$. Each original pair subgroup contains
$Z$, so its order is twice the corresponding projective order. The subgroup generated by
all three has order divisible by $\operatorname{lcm}(8,12,20)=120$, and is
therefore $H$. Hence the three matrices are independent and $\mup(H)=3$.
Their three omission subgroups intersect in $Z$, and each is essential.
Indeed, the first two contain $\langle R_3\rangle$, of order 4, and their
intersection has order dividing $\gcd(8,12)=4$, so it is exactly
$\langle R_3\rangle$. The third cannot contain $R_3$. Thus they also give $b(H)\ge3$. The upper bound
$b(H)\le\mup(H)$ finishes the proof.
\end{proof}

\section{A small affine lemma}
\begin{lemma}\label{lem:line}
Let $p$ be an odd prime and let $C$ be a cyclic 2-subgroup of $\F_p^\times$.
For the faithful scalar semidirect product $L=\F_p\rtimes C$,
$\mup(L)\le2$.
\end{lemma}
\begin{proof}
Put $P=(\F_p,+)$. A subgroup $J$ with $J\cap P=1$ injects into $C$, so is a
cyclic 2-group and has independence number at most 1. Otherwise $J$ contains
$P$, and is $P\rtimes D$ for some $D\le C$. The case $D=1$ is cyclic of
prime order. In the other cases, every generating set for $J$ contains an
element $x$ whose projection generates $D$, because a cyclic 2-group has a
unique maximal subgroup. If $x=(a,\lambda)$ and $|D|=m$, then
\[
 x^m=\left(a\sum_{j=0}^{m-1}\lambda^j,1\right)=1,
\]
so $\langle x\rangle$ is a complement of order $m$. Some other generator
$y$ is outside this complement. The subgroup $\langle x,y\rangle$ projects
onto $D$; if its intersection with $P$ were trivial it would have order $m$
and would equal $\langle x\rangle$, a contradiction. Thus it contains $P$
and equals $J$. Every independent generating set of every subgroup has size
at most 2.
\end{proof}

\begin{lemma}\label{lem:stabilizers}
Every line stabilizer in $H\le\SL_2(29)$ is cyclic of order dividing 4.
Consequently $H$ acts irreducibly on $V=\F_{29}^2$.
\end{lemma}
\begin{proof}
For a line $W$, the scalar action on $W$ defines a homomorphism
$H_W\to\F_{29}^\times$. In a basis adapted to $W$, an element in its kernel
has the form $\left(\begin{smallmatrix}1&u\\0&1\end{smallmatrix}\right)$,
because its determinant is one. If $u\ne0$, this element has order 29,
which cannot divide $|H|=120$. Thus the homomorphism is injective, and
\[
                 |H_W|\mid\gcd(120,28)=4.
\]
Since $|H|>4$, the whole group cannot stabilize a line.
\end{proof}

\section{The affine counterexample and the all-actions upper bound}
Set
\[
 G=V\rtimes H,\qquad V=\F_{29}^2,
\]
with multiplication
\[
              (v,h)(w,k)=(v+hw,hk).
\]
The group order is $29^2\cdot120=100920$; write $\pi:G\to H$ for projection.
This description specifies the action, not just the abstract factors of a
semidirect product.

\begin{lemma}\label{lem:properV}
If $K\le G$ does not contain $V$, then $\bfth(K)\le2$. More precisely, if
$K\cap V=1$, then $\mup(K)\le3$ and $\bfth(K)\le2$; if
$\dim_{\F_{29}}(K\cap V)=1$, then $\mup(K)\le2$.
\end{lemma}
\begin{proof}
Every additive subgroup of $V$ is an $\F_{29}$-subspace, since 29 is prime.
If $K\cap V=1$, then $\pi$ embeds $K$ in $H$. Both $\mup$ and $\bfth$ are
monotone for subgroups: for the latter, use the same subgroup family with
trivial intersection in the larger group. Apply Lemma~\ref{lem:H}.

Otherwise $W=K\cap V$ is a line, normal in $K$, so $\pi(K)\le H_W$ is cyclic
of order dividing 4. A Sylow 2-subgroup of $K$ maps isomorphically onto
$\pi(K)$ and complements $W$. Its action on $W$ is faithful by
Lemma~\ref{lem:stabilizers}. Thus $K$ is isomorphic to a group covered by
Lemma~\ref{lem:line}, which gives $\mup(K)\le2$ and hence $\bfth(K)\le2$.
These are all possibilities for a subgroup not containing $V$.
\end{proof}

\newpage
\begin{proposition}\label{prop:faithful}
$\bfth(G)\le3$.
\end{proposition}
\begin{proof}
Let $K_1,\ldots,K_t$ be any meet-irredundant family with total intersection
$1$. Some $K_i$ does not contain $V$, since otherwise the intersection would
contain $V$. Lemmas~\ref{lem:restriction} and \ref{lem:properV} give
\[
                   t-1\le\bfth(K_i)\le2.
\]
No transitivity, maximality, or conjugacy assumption on the subgroups was made.
\end{proof}

\begin{proposition}\label{prop:allactions}
$b(G)\le3$, including all nonfaithful actions.
\end{proposition}
\begin{proof}
Every nontrivial normal subgroup $N\lhd G$ contains $V$. To prove this,
irreducibility gives $N\cap V=1$ or $V$. In the first case
$[N,V]\le N\cap V=1$, so $N\le C_G(V)$. The matrix action is faithful, hence
$C_G(V)=V$, forcing $N=1$. This contradicts the choice of $N$.

It follows that every nonfaithful action of $G$ factors through $G/V=H$,
and its minimal bases have size at most $b(H)=3$. Faithful actions are
covered by Proposition~\ref{prop:faithful}. This proves the stated bound
for every representation.
\end{proof}

\section{Explicit lower bounds and the exact independence number}
Let $e_1=(1,0)^\mathsf{T}$ and $e_2=(0,1)^\mathsf{T}$. With the matrices in
\eqref{eq:R}, define
\[
 S=\{(e_1,I),(0,R_1),(0,R_2),(0,R_3)\}\subseteq G.
\]
The three linear elements project to the independent triple in $H$. Thus no
one of them is generated by the other three elements of $S$, since its
projection would then be generated by the other two matrices. The translation
$(e_1,I)$ is not in the zero-translation complement. Therefore $S$ is
independent and $\mup(G)\ge4$.

In fact $S$ generates $G$: it contains generators for $H$, and the orbit of
a nonzero vector under an irreducible group spans $V$. Its omission subgroups
have the following explicitly verified orders:
\begin{center}
\begin{tabular}{lc}
\toprule
Omitted element & Order of the subgroup generated by the others\\
\midrule
$(e_1,I)$ & $120$\\
$(0,R_1)$ & $841\cdot8=6728$\\
$(0,R_2)$ & $841\cdot12=10092$\\
$(0,R_3)$ & $841\cdot20=16820$\\
\bottomrule
\end{tabular}
\end{center}
The latter three formulas also follow structurally: each corresponding pair
subgroup of $H$ has order greater than 4, so cannot stabilize a line and is
irreducible on $V$.

\begin{proposition}\label{prop:mu}
$\mup(G)=4$.
\end{proposition}
\begin{proof}
It remains to give the upper bound. By Lemma~\ref{lem:breadth}, consider any
meet-irredundant family $K_1,\ldots,K_t$, with no assumption on its total
intersection. If all members contain $V$, quotienting by the genuinely normal
subgroup $V$ gives an irredundant family in $H$, so $t\le3$.
Otherwise choose $K_i\not\supseteq V$. If $K_i\cap V=1$, then
$t-1\le\mup(K_i)\le3$. If $K_i\cap V$ is a line, then
$t-1\le\mup(K_i)\le2$. In either case $t\le4$.
\end{proof}

There is also a faithful minimal base of size 3. Define subgroups
\[
\begin{split}
 L_1&=\langle(e_1,I),(0,-I)\rangle,\\
 L_2&=\langle(e_2,I),(0,-I)\rangle,\\
 L_3&=\langle(e_1+e_2,I),(2e_1,-I)\rangle.
\end{split}
\]
Each has order 58. Their pairwise intersections are
\[
\begin{split}
 L_1\cap L_2&=\{1,(0,-I)\},\\
 L_1\cap L_3&=\{1,(2e_1,-I)\},\\
 L_2\cap L_3&=\{1,(-2e_2,-I)\}.
\end{split}
\]
Their total intersection is $1$. The three identity cosets in the disjoint
union $G/L_1\sqcup G/L_2\sqcup G/L_3$ therefore form a faithful minimal base
of size 3. This particular action has degree $3(100920/58)=5220$.

\begin{theorem}[Counterexample]
For the explicitly specified group
\[
 G=\F_{29}^2\rtimes
 \left\langle
 \begin{pmatrix}0&28\\1&0\end{pmatrix},
 \begin{pmatrix}2&7\\12&28\end{pmatrix}
 \right\rangle,
\]
one has
\[
             |G|=100920,
 \qquad      \bfth(G)=b(G)=3<4=\mup(G).
\]
This gives a negative answer to Kourovka Notebook Problem 16.45.
\end{theorem}

\section{The normality obstruction is explicit}
Let $M_i=\langle S\setminus\{s_i\}\rangle$ for the four elements of $S$ in the
order displayed above. Then
\[
              \bigcap_{i=1}^4M_i=Z_0,
       \qquad Z_0=\{(0,I),(0,-I)\}.
\]
Indeed, $M_1=H$ and the other three are the preimages of the three pair
subgroups in Lemma~\ref{lem:H}, whose intersection is $Z$. The subgroup
$Z_0$ is not normal in $G$, because
\[
 (e_1,I)(0,-I)(-e_1,I)=(2e_1,-I)\notin Z_0.
\]
Thus this independent 4-set gives an irredundant intersection family, and a
Boolean meet-semilattice of rank 4, with nonnormal bottom. The theorem proves
that no rank-four embedding can have normal bottom.

Taking normal cores does not repair this example. Simplicity of $H/Z\cong A_5$
implies that the three proper pair subgroups of $H$ have core $Z$. Since
$M_2,M_3,M_4$ contain $V$, their cores in $G$ are $V\rtimes Z$.
The complement $M_1=H$ has trivial core, since every nontrivial normal
subgroup of $G$ contains $V$, while $H\cap V=1$. Hence the four cores are
\[
                   1,\quad V\rtimes Z,\quad V\rtimes Z,\quad V\rtimes Z.
\]
The first member alone gives their total intersection. All four-way
essentiality has been lost. No quotient by the nonnormal subgroup $Z_0$ is
used anywhere in the argument.

\section{Exact certification and a boundary case}
The supplied Python and C++ verifiers are separate implementations and use
exact integers only. Python uses
only its standard library. The second implementation is C++17 and reads no
Python-produced group table. Both reconstruct the complement and enumerate
\emph{all 76 actual subgroups}, not just conjugacy representatives. Their
subgroup sets were compared as sets of actual matrices and agree.

Coverage of the subgroup enumeration follows by adjoining every element
outside every subgroup already found, starting with the identity subgroup.
For any subgroup $K$, a generating sequence for $K$ is a path in this search.
The two verifiers check all $\binom{75}{4}=1215450$ four-element families of
proper subgroups of $H$ and find none meet-irredundant. Every subfamily of an
irredundant family is irredundant, so this excludes all larger ranks as well.
They then check all $\binom{75}{3}=67525$ three-element families and find none
both faithful and minimal. Together with the four-family check, this
independently certifies $\mup(H)=3$ and $\bfth(H)=2$ using the explicit lower
witnesses. The mathematical proofs above do not require trusting an optimizer.

They also verify the independent affine 4-set, generate all 100920 elements
of $G$ without storing its full multiplication table, check the omission
orders and the nonnormal intersection, and verify the faithful 3-family.
The Python verifier explicitly checks all 30 line stabilizers, all of which
are in fact cyclic of order 4, and includes the following adversarial test.

The characteristic restriction cannot be discarded. Over $\F_{11}$, use
\[
 A_{11}=\begin{pmatrix}0&10\\1&0\end{pmatrix},\quad
 B_{11}=\begin{pmatrix}0&2\\5&1\end{pmatrix},\quad
 C=\begin{pmatrix}6&10\\0&2\end{pmatrix}.
\]
The first two generate a complement of order 120; $C$ has order 10 and
stabilizes the distinct lines $W_1=\langle(1,0)\rangle$ and
$W_2=\langle(1,4)\rangle$. In $\F_{11}^2\rtimes\langle A_{11},B_{11}\rangle$,
the four subgroups
\[
 W_1\rtimes\langle C\rangle,\quad
 W_2\rtimes\langle C\rangle,\quad
 V\rtimes\langle C^5\rangle,\quad
 V\rtimes\langle C^2\rangle
\]
have total intersection $1$. Their deletion intersections have respective
orders $11,11,5,2$, so they give a faithful minimal base of size 4. This
explicitly falsifies the tempting but incorrect claim that every such
characteristic gives the counterexample. The equality
$\gcd(120,28)=4$, and the resulting line-subgroup bound, are essential.

Earlier exploratory runs used integer-programming models on complete small
subgroup lattices. Their code and logs are retained for provenance, but no
floating-point MILP upper bound is a premise of this proof. The exact
certificate and the structural all-actions argument replace those exploratory
results. The accompanying audit directory records a separate structural review,
a newly written exact checker, and fresh replays of the supplied verification
programs. These are AI-assisted checks, not external human peer review or a
proof-assistant formalization. The structural lemmas and all-actions argument
remain mathematical proof obligations for a reader; finite computations
certify the explicit inputs and witnesses.

\newpage
\section*{Author and reproducibility statement}
Alec Kriebel is the author. This publication revision and its verification
were prepared with OpenAI Codex, including separate AI agents for structural,
computational, and source audits. No external researcher was contacted during
this review. The manuscript is an unrefereed preprint.

The source, exact certificates, independent checks, and audit reports are
available with this paper at
\url{https://aleckriebel.github.io/Math/papers/kourovka-16-45/}.
From the source directory, \texttt{make verify} runs the two supplied exact
verifiers and compares their complete subgroup sets;
\texttt{python3 audit\_2026\_09\_21/independent\_check.py} runs the additional
checker. Python uses only the standard library, and the second supplied
implementation requires a C++17 compiler. Neither the mathematical proof nor
these final checks depend on the exploratory optimization software.

\begin{thebibliography}{9}
\bibitem{notebook}
E.~I. Khukhro and V.~D. Mazurov (eds.),
\emph{Unsolved Problems in Group Theory: The Kourovka Notebook}, 21st edition,
September 2026 update, Problem 16.45, printed p.~100.
\url{https://kourovkanotebookorg.wordpress.com/wp-content/uploads/2026/09/21tkt.pdf}.
Separate update:
\url{https://kourovkanotebookorg.wordpress.com/wp-content/uploads/2026/09/21upd.pdf}.

\bibitem{blog}
P.~J. Cameron, \emph{The symmetric group, 7}, 22 July 2010.
\url{https://cameroncounts.wordpress.com/2010/07/22/the-symmetric-group-7/}.

\bibitem{paper}
P.~J. Cameron, \emph{Some measures of finite groups related to permutation
bases}, arXiv:1408.0968, 2014, especially Propositions 3.1--3.2 and
Corollary 3.3. \url{https://arxiv.org/abs/1408.0968}.

\bibitem{cameron2024}
P.~J. Cameron, \emph{Independence and bases: Theme and variations},
Model Theory \textbf{3} (2024), no.~2, 417--431.
\url{https://doi.org/10.2140/mt.2024.3.417}.

\bibitem{blog2014}
P.~J. Cameron, \emph{Groups, lattices and bases}, 6 August 2014.
\url{https://cameroncounts.wordpress.com/2014/08/06/groups-lattices-and-bases/}.
\end{thebibliography}
\end{document}
